\documentclass[10pt,a4paper]{amsart}
\usepackage[margin=19mm]{geometry}
\usepackage{amsmath,amssymb,mathtools}
\usepackage[scr=rsfs,cal=euler]{mathalfa}
\usepackage{xfrac}
\usepackage[unicode,hidelinks,bookmarks=false]{hyperref}
\newcommand{\ie}{i.e.\ }
\newcommand{\cf}{cf.\ }
\newcommand{\resp}{respectively\ }
\newcommand{\Subsection}{\subsection}
\providecommand{\nobreakdash}{\nobreak\hbox{-}}
\setlength{\emergencystretch}{2em}
\multlinegap0pt
\usepackage{lmodern}
\newtheorem{theorem}{Theorem}[section]
\newtheorem{corollary}[theorem]{Corollary}
\newtheorem*{remark}{Remark}
\newtheorem{example}[theorem]{Example}
\newcommand{\Rad}{\operatorname{Rad}}
\newcommand{\ltr}{\operatorname{ltr}}
\begin{document}
\title[Parallel Radical and Screen Distributions on Lightlike Hyper]{Parallel Radical and Screen Distributions on Lightlike Hypersurfaces of Indefinite Statistical Manifolds}
\author{Burçin Doğan}
\date{}
\maketitle
\noindent\emph{Source and licence.} Parallel Radical and Screen Distributions on Lightlike Hypersurfaces of Indefinite Statistical Manifolds. Version arXiv:2609.17931v1 (CC BY 4.0). This material is distributed under the Creative Commons Attribution 4.0 International licence, http://creativecommons.org/licenses/by/4.0/, as recorded in the arXiv OAI licence metadata for this exact version and shown on the abstract page https://arxiv.org/abs/2609.17931v1. Full rights notice: licenses/arxiv-2609.17931-CC-BY-4.0.txt.\par\smallskip
\noindent\textit{Editorial scope.} Counted unit: the original \texttt{theorem} environment \texttt{thm:K-splitting} at source lines 568--590 together with its complete proof at lines 592--611; the delivered document keeps source lines 181--612 so that every referenced label and every citation resolves inside it. Native editable source is the paper's own concatenated TeX; the AMS wrapper is editorial formatting only. No publisher class, upstream script or unrelated artwork is redistributed.
\begin{equation}\label{eq:duality}
X\overline{g}(Y,Z)
=
\overline{g}(\overline{\nabla}_{X}Y,Z)
+
\overline{g}(Y,\overline{\nabla}^{*}_{X}Z).
\end{equation}
If $(\overline{\nabla},\overline{g})$ is statistical, then
$\overline{\nabla}^{*}$ is also torsion-free and
$(\overline{\nabla}^{*},\overline{g})$ is a statistical structure.
Moreover, the Levi--Civita connection $\overline{\nabla}^{0}$ of
$\overline{g}$ satisfies
\begin{equation}\label{eq:levicivita-mean}
\overline{\nabla}^{0}
=
\frac{1}{2}
\left(
\overline{\nabla}
+
\overline{\nabla}^{*}
\right).
\end{equation}
We refer to
$(\overline{M},\overline{g},\overline{\nabla},
\overline{\nabla}^{*})$
as an indefinite statistical manifold
\cite{Amari1985,Lauritzen1987,Vos1989,Furuhata2009}.

Let $M$ be a lightlike hypersurface of
$(\overline{M},\overline{g})$, and denote by $g$ the degenerate metric
induced on $M$. Since $M$ has codimension one, its radical distribution
is the rank-one bundle
\begin{equation}\label{eq:radical}
\Rad(TM)=TM\cap TM^{\perp}=TM^{\perp}.
\end{equation}
A screen distribution $S(TM)$ is a non-degenerate complementary
distribution to $\Rad(TM)$ in $TM$. Thus,
\begin{equation}\label{eq:tangent-decomposition}
TM=S(TM)\mathbin{\perp}\Rad(TM).
\end{equation}
Choose a nowhere-zero local section $\xi$ of $\Rad(TM)$. Corresponding
to the choice of $S(TM)$, there exists locally a unique lightlike
transversal vector field $N$ satisfying
\begin{equation}\label{eq:null-frame}
\overline{g}(N,N)=0,\qquad
\overline{g}(N,\xi)=1,\qquad
\overline{g}(N,W)=0
\end{equation}
for every $W\in\Gamma(S(TM))$. The line bundle generated by $N$ is
denoted by $\ltr(TM)$. Consequently,
\begin{equation}\label{eq:ambient-decomposition}
T\overline{M}|_{M}
=
S(TM)\mathbin{\perp}
\bigl(\Rad(TM)\oplus\ltr(TM)\bigr)
=
TM\oplus\ltr(TM).
\end{equation}
These constructions are standard in lightlike hypersurface geometry
\cite{DuggalBejancu1996,DuggalJin2007,DuggalSahin2010}.

Let $P$ denote the projection of $TM$ onto $S(TM)$ with respect to
\eqref{eq:tangent-decomposition}, and put
\begin{equation}\label{eq:eta}
\eta(X)=\overline{g}(X,N),\qquad X\in\Gamma(TM).
\end{equation}
Then every tangent vector field admits the decomposition
\begin{equation}\label{eq:projection}
X=PX+\eta(X)\xi.
\end{equation}

For $X,Y\in\Gamma(TM)$, the Gauss formulas associated with the dual
ambient connections are written as
\begin{align}
\overline{\nabla}_{X}Y
&=
\nabla_{X}Y+B(X,Y)N,
\label{eq:gauss}\\
\overline{\nabla}^{*}_{X}Y
&=
\nabla^{*}_{X}Y+B^{*}(X,Y)N,
\label{eq:gauss-star}
\end{align}
where $\nabla$ and $\nabla^{*}$ are the induced torsion-free
connections on $M$, and $B$ and $B^{*}$ are the corresponding second
fundamental forms. Since the ambient connections are torsion-free,
$B$ and $B^{*}$ are symmetric. Following the convention used for lightlike hypersurfaces of statistical
manifolds in \cite{BahadirTripathi2023}, the Weingarten formulas are
\begin{align}
\overline{\nabla}_{X}N
&=
-A^{*}_{N}X+\tau^{*}(X)N,
\label{eq:weingarten}\\
\overline{\nabla}^{*}_{X}N
&=
-A_{N}X+\tau(X)N,
\label{eq:weingarten-star}
\end{align}
where $A^{*}_{N}$ and $A_{N}$ are the corresponding Weingarten mappings
and $\tau^{*},\tau$ are one-forms on $M$.

The induced connections can be decomposed further along the screen
distribution. For $X\in\Gamma(TM)$ and $W\in\Gamma(S(TM))$, write
\begin{align}
\nabla_{X}W
&=
\nabla^{S}_{X}W+C(X,W)\xi,
\label{eq:screen-gauss}\\
\nabla^{*}_{X}W
&=
\nabla^{*S}_{X}W+C^{*}(X,W)\xi,
\label{eq:screen-gauss-star}
\end{align}
where $\nabla^{S}$ and $\nabla^{*S}$ are the induced connections on $S(TM)$,
while $C,C^{*}$ are the corresponding local screen fundamental forms.
For the radical vector field $\xi$, the tangential parts of the
derivatives are expressed as
\begin{align}
\nabla_{X}\xi
&=
-A_{\xi}X-\tau(X)\xi,
\label{eq:xi-derivative}\\
\nabla^{*}_{X}\xi
&=
-A^{*}_{\xi}X-\tau^{*}(X)\xi,
\label{eq:xi-derivative-star}
\end{align}
where $A_{\xi}$ and $A^{*}_{\xi}$ take values in $S(TM)$ and are the
corresponding screen shape operators.

The duality relation \eqref{eq:duality} yields, for
$X,Y,Z\in\Gamma(TM)$,
\begin{align}
Xg(Y,Z)
={}&
g(\nabla_{X}Y,Z)
+
g(Y,\nabla^{*}_{X}Z) \nonumber\\
&\quad
+B(X,Y)\eta(Z)
+
B^{*}(X,Z)\eta(Y).
\label{eq:induced-duality}
\end{align}
Thus, in contrast with the non-degenerate statistical submanifold case,
the induced connections $\nabla$ and $\nabla^{*}$ need not form a dual
pair with respect to the degenerate metric $g$. On the screen
distribution, however, \eqref{eq:induced-duality} reduces to
\begin{equation}\label{eq:screen-duality}
Xg(U,V)
=
g(\nabla^{S}_{X}U,V)
+
g(U,\nabla^{*S}_{X}V),
\qquad
U,V\in\Gamma(S(TM)).
\end{equation}
Hence $\nabla^{S}$ and $\nabla^{*S}$ are dual with respect to the
non-degenerate metric induced on $S(TM)$.

The second fundamental forms and shape operators satisfy the relations
\begin{align}
B(X,W)
&=
g(A^{*}_{\xi}X,W),
&
B^{*}(X,W)
&=
g(A_{\xi}X,W),
\label{eq:B-shape}\\
C(X,W)
&=
g(A_{N}X,W),
&
C^{*}(X,W)
&=
g(A^{*}_{N}X,W),
\label{eq:C-shape}
\end{align}
for $X\in\Gamma(TM)$ and $W\in\Gamma(S(TM))$. In addition,
\begin{equation}\label{eq:B-xi}
B(X,\xi)+B^{*}(X,\xi)=0.
\end{equation}
Related formulas for the induced geometric objects can be found in
\cite{JainSinghKumar2020,KazemiSalahvarzi2020,RaniKaur2021,
BahadirTripathi2023}.

For later use, define the mean induced connection and the
difference tensor by
\begin{equation}\label{eq:mean-difference}
\nabla^{0}=\frac{1}{2}\left(\nabla+\nabla^{*}\right),
\qquad
K(X,Y)=\frac{1}{2}\left(\nabla_{X}Y-\nabla^{*}_{X}Y\right).
\end{equation}
Then
\begin{equation}\label{eq:connection-decomposition}
\nabla_{X}Y=\nabla^{0}_{X}Y+K(X,Y),
\qquad
\nabla^{*}_{X}Y=\nabla^{0}_{X}Y-K(X,Y).
\end{equation}
Since $\nabla$ and $\nabla^{*}$ are torsion-free, $K$ is symmetric in its
two entries. The connection $\nabla^{0}$ is the induced connection obtained from the
Levi--Civita connection $\overline{\nabla}^{0}$ along $M$. If
$B^{0}=\frac12(B+B^{*})$, then
\begin{equation}\label{eq:nabla0-g}
(\nabla^{0}_{X}g)(Y,Z)
=
B^{0}(X,Y)\eta(Z)+B^{0}(X,Z)\eta(Y),
\qquad B^{0}(X,\xi)=0.
\end{equation}
Hence $\nabla^{0}$ is a metric connection if and only if $B^{0}=0$.
Since the induced metric is degenerate, $\nabla^{0}$ is not an intrinsic
Levi--Civita connection of $(M,g)$.

% ================================================================
\section{Parallel Radical and Screen Distributions}
\label{sec:parallel}
% ================================================================

Let $M$ be a lightlike hypersurface of an indefinite statistical manifold
$(\overline{M},\overline{g},\overline{\nabla},
\overline{\nabla}^{*})$. A distribution $\mathcal D$ on $M$ is said to
be parallel with respect to $\nabla$ if
$\nabla_XY\in\Gamma(\mathcal D)$ for all
$X\in\Gamma(TM)$ and $Y\in\Gamma(\mathcal D)$; the definition for
$\nabla^{*}$ is analogous. Related conditions for radical and screen distributions with respect to
individual induced connections appear in
\cite{JainSinghKumar2020,RaniKaur2021,BahadirTripathi2023}. Here we keep
the two induced connections distinct and focus on the conditions that hold for
both of them, comparing these common conditions through $\nabla^{0}$ and $K$.

\begin{theorem}\label{thm:radical-both}
Let $M$ be a lightlike hypersurface of an indefinite statistical manifold
$(\overline{M},\overline{g},\overline{\nabla},\overline{\nabla}^{*})$.
Denote by $\nabla$ and $\nabla^{*}$ the connections on $M$ induced by
$\overline{\nabla}$ and $\overline{\nabla}^{*}$, respectively. The following
assertions are equivalent:
\begin{enumerate}
\item $\Rad(TM)$ is parallel with respect to both $\nabla$ and
$\nabla^{*}$;
\item $A_{\xi}=A_{\xi}^{*}=0$;
\item $\nabla^{0}$ is a metric connection and
$K(X,\xi)\in\Gamma(\Rad(TM))$ for all $X\in\Gamma(TM)$;
\item $B(X,W)=B^{*}(X,W)=0$ for all $X\in\Gamma(TM)$ and
$W\in\Gamma(S(TM))$.
\end{enumerate}
\end{theorem}

\begin{proof}
Equations~\eqref{eq:xi-derivative} and
\eqref{eq:xi-derivative-star} give (1)$\Leftrightarrow$(2). For $W\in\Gamma(S(TM))$, \eqref{eq:B-shape} shows that
$2B^{0}(X,W)$ equals $g((A_{\xi}+A_{\xi}^{*})X,W)$. Together with
$B^{0}(X,\xi)=0$ in \eqref{eq:nabla0-g} and the fact that the metric on
$S(TM)$ is non-degenerate, this shows that $B^{0}=0$ if and only if
$A_{\xi}+A_{\xi}^{*}=0$. Thus, by \eqref{eq:nabla0-g},
$\nabla^{0}$ is a metric connection if and only if
$A_{\xi}+A_{\xi}^{*}=0$. Moreover,
\begin{equation}\label{eq:K-xi}
K(X,\xi)
=-\frac12(A_{\xi}-A_{\xi}^{*})X
-\frac12\{\tau(X)-\tau^{*}(X)\}\xi.
\end{equation}
Hence (2)$\Leftrightarrow$(3). Finally, \eqref{eq:B-shape} gives
$A_{\xi}=0$ if and only if $B^{*}(X,W)=0$, and
$A_{\xi}^{*}=0$ if and only if $B(X,W)=0$. Thus
(2)$\Leftrightarrow$(4).
\end{proof}

\begin{corollary}\label{cor:radical-totally-geodesic}
Let $M$ be a lightlike hypersurface of
$(\overline{M},\overline{g},\overline{\nabla},\overline{\nabla}^{*})$ with
induced connections $\nabla$ and $\nabla^{*}$. If $\Rad(TM)$ is parallel
with respect to both induced connections, then
\begin{equation}\label{eq:B-radical-form}
B(X,Y)=\eta(X)\eta(Y)B(\xi,\xi),
\qquad
B^{*}(X,Y)=-\eta(X)\eta(Y)B(\xi,\xi).
\end{equation}
Consequently, $M$ is totally geodesic with respect to
$\overline{\nabla}$ if and only if it is totally geodesic with respect
to $\overline{\nabla}^{*}$; either condition is equivalent to
$B(\xi,\xi)=0$.
\end{corollary}

\begin{proof}
Write $X=PX+\eta(X)\xi$ and $Y=PY+\eta(Y)\xi$. By
Theorem~\ref{thm:radical-both}, $B$ and $B^{*}$ vanish whenever one
entry lies in $S(TM)$. Their symmetry then leaves only the
radical--radical terms, and \eqref{eq:B-xi} gives
\eqref{eq:B-radical-form}. The last assertion
follows from the definitions of a totally geodesic lightlike hypersurface
with respect to the two ambient connections.
\end{proof}

\begin{remark}
The identities in \eqref{eq:B-shape} are used with the second entry in
$S(TM)$. Thus the vanishing of $A_{\xi}$ and $A_{\xi}^{*}$ eliminates all
components of $B$ and $B^{*}$ with one entry in $S(TM)$, but does not
by itself eliminate the remaining radical--radical component. This is why the additional condition
$B(\xi,\xi)=0$ appears in Corollary~\ref{cor:radical-totally-geodesic}.
\end{remark}

The separate conditions involving $C$ and $C^{*}$ are closely related to
those considered in \cite{RaniKaur2021,BahadirTripathi2023}. The next result keeps the
two induced connections distinct and records their common condition.

\begin{theorem}\label{thm:screen-both}
Consider a lightlike hypersurface $M$ of the indefinite statistical manifold
$(\overline{M},\overline{g},\overline{\nabla},\overline{\nabla}^{*})$, and
write $\nabla$ and $\nabla^{*}$ for the induced connections. The following
assertions are equivalent:
\begin{enumerate}
\item $S(TM)$ is parallel with respect to both $\nabla$ and
$\nabla^{*}$;
\item $C=C^{*}=0$ on $TM\times S(TM)$;
\item $A_NX,A_N^{*}X\in\Gamma(\Rad(TM))$ for all
$X\in\Gamma(TM)$;
\item $S(TM)$ is parallel with respect to $\nabla^{0}$ and
$K(X,W)\in\Gamma(S(TM))$ for all $X\in\Gamma(TM)$ and
$W\in\Gamma(S(TM))$.
\end{enumerate}
\end{theorem}

\begin{proof}
Equations~\eqref{eq:screen-gauss}--\eqref{eq:screen-gauss-star} give
(1)$\Leftrightarrow$(2). By \eqref{eq:C-shape}, and since the metric induced on $S(TM)$ is
non-degenerate, (2)$\Leftrightarrow$(3). Finally,
\begin{align}
\nabla_X^{0}W
&=\frac12(\nabla_X^{S}W+\nabla_X^{*S}W)
 +\frac12\{C(X,W)+C^{*}(X,W)\}\xi,\label{eq:nabla0-screen}\\
K(X,W)
&=\frac12(\nabla_X^{S}W-\nabla_X^{*S}W)
 +\frac12\{C(X,W)-C^{*}(X,W)\}\xi.\label{eq:K-screen}
\end{align}
Thus (4) holds if and only if $C+C^{*}=0$ and $C-C^{*}=0$, which is
precisely (2).
\end{proof}

We now consider the two distributions simultaneously. The screen
projection gives a concise formulation of the resulting condition.

\begin{theorem}\label{thm:both-distributions}
Let $M$, with induced connections $\nabla$ and $\nabla^{*}$, be a lightlike
hypersurface of the indefinite statistical manifold
$(\overline{M},\overline{g},\overline{\nabla},\overline{\nabla}^{*})$.
The following assertions are equivalent:
\begin{enumerate}
\item $\Rad(TM)$ and $S(TM)$ are parallel with respect to both
$\nabla$ and $\nabla^{*}$;
\item the screen projection $P$ is parallel with respect to both
$\nabla$ and $\nabla^{*}$;
\item $A_{\xi}=A_{\xi}^{*}=0$ and $C=C^{*}=0$ on
$TM\times S(TM)$;
\item $P$ is parallel with respect to $\nabla^{0}$ and
\begin{equation}\label{eq:K-commutes-P}
K(X,PY)=P K(X,Y)
\end{equation}
for all $X,Y\in\Gamma(TM)$.
\end{enumerate}
\end{theorem}

\begin{proof}
Using \eqref{eq:projection}, \eqref{eq:screen-gauss}, and
\eqref{eq:xi-derivative}, we obtain
\begin{equation}\label{eq:nabla-P}
(\nabla_XP)Y=C(X,PY)\xi+\eta(Y)A_{\xi}X.
\end{equation}
Similarly,
\begin{equation}\label{eq:nabla-star-P}
(\nabla_X^{*}P)Y=C^{*}(X,PY)\xi+\eta(Y)A_{\xi}^{*}X.
\end{equation}
By Theorems~\ref{thm:radical-both} and \ref{thm:screen-both},
(1) is equivalent to (3). On the other hand, taking first
$Y\in\Gamma(S(TM))$ and then $Y=\xi$ in
\eqref{eq:nabla-P}--\eqref{eq:nabla-star-P} shows that (2) is
equivalent to (3). From \eqref{eq:connection-decomposition},
\begin{align*}
(\nabla_XP)Y
&=(\nabla_X^{0}P)Y+K(X,PY)-P K(X,Y),\\
(\nabla_X^{*}P)Y
&=(\nabla_X^{0}P)Y-K(X,PY)+P K(X,Y).
\end{align*}
They vanish simultaneously if and only if $\nabla^{0}P=0$ and
\eqref{eq:K-commutes-P} holds. Hence (2)$\Leftrightarrow$(4).
\end{proof}


\begin{theorem}\label{thm:K-splitting}
Let $M$ be a lightlike hypersurface of
$(\overline{M},\overline{g},\overline{\nabla},\overline{\nabla}^{*})$.
Suppose that the connections $\nabla$ and $\nabla^{*}$ induced by
$\overline{\nabla}$ and $\overline{\nabla}^{*}$ preserve both
$\Rad(TM)$ and $S(TM)$. Then the mixed radical--screen part of the
difference tensor vanishes:
\begin{equation}\label{eq:K-mixed-zero}
K(W,\xi)=K(\xi,W)=0,
\qquad W\in\Gamma(S(TM)).
\end{equation}
Moreover, for all $X,Y\in\Gamma(TM)$,
\begin{equation}\label{eq:K-block-decomposition}
K(X,Y)=K(PX,PY)+\eta(X)\eta(Y)K(\xi,\xi),
\end{equation}
where $K(PX,PY)\in\Gamma(S(TM))$ and
$K(\xi,\xi)\in\Gamma(\Rad(TM))$. In addition,
\begin{equation}\label{eq:tau-difference-screen}
(\tau-\tau^{*})|_{S(TM)}=0,
\qquad
\tau-\tau^{*}=\{\tau(\xi)-\tau^{*}(\xi)\}\eta.
\end{equation}
\end{theorem}

\begin{proof}
By Theorem~\ref{thm:both-distributions},
$K(X,PY)=PK(X,Y)$. For $W\in\Gamma(S(TM))$, taking $X=\xi$ and
$Y=W$ gives $K(\xi,W)=PK(\xi,W)$, so $K(\xi,W)$ is screen-valued.
Taking instead $X=W$ and $Y=\xi$ gives $PK(W,\xi)=0$. Since $K$ is
symmetric, $K(\xi,W)=K(W,\xi)$ is therefore both screen-valued and
radical-valued, and hence it vanishes. This proves
\eqref{eq:K-mixed-zero}.

Expanding $X=PX+\eta(X)\xi$ and $Y=PY+\eta(Y)\xi$ and using
\eqref{eq:K-mixed-zero} gives \eqref{eq:K-block-decomposition}. The first
term is screen-valued by \eqref{eq:K-commutes-P}, whereas setting
$Y=\xi$ in that relation gives $PK(X,\xi)=0$, and in particular
$K(\xi,\xi)\in\Gamma(\Rad(TM))$. Finally, Theorem~\ref{thm:radical-both}
and \eqref{eq:K-xi} imply
$K(W,\xi)=-\frac12\{\tau(W)-\tau^{*}(W)\}\xi$ for
$W\in\Gamma(S(TM))$. Together with \eqref{eq:K-mixed-zero}, this proves
the first identity in \eqref{eq:tau-difference-screen}; the second follows
from $X=PX+\eta(X)\xi$ and $\eta(\xi)=1$.
\end{proof}

\begin{thebibliography}{99}
\bibitem[Amari1985]{Amari1985} S.-I. Amari, \textit{Differential-geometrical methods in statistics}, Lecture Notes in Statistics, vol.~28, Springer-Verlag, New York, 1985.
\bibitem[BahadirTripathi2023]{BahadirTripathi2023} O. Bahad{\i}r and M. M. Tripathi, Geometry of lightlike hypersurfaces of a statistical manifold, \textit{WSEAS Trans. Math.} \textbf{22} (2023), 466--474, \url{https://doi.org/10.37394/23206.2023.22.52}.
\bibitem[DuggalBejancu1996]{DuggalBejancu1996} K. L. Duggal and A. Bejancu, \textit{Lightlike submanifolds of semi-Riemannian manifolds and applications}, Mathematics and Its Applications, vol.~364, Kluwer Academic Publishers, Dordrecht, 1996.
\bibitem[DuggalJin2007]{DuggalJin2007} K. L. Duggal and D. H. Jin, \textit{Null curves and hypersurfaces of semi-Riemannian manifolds}, World Scientific Publishing Co., Hackensack, NJ, 2007.
\bibitem[DuggalSahin2010]{DuggalSahin2010} K. L. Duggal and B. \c{S}ahin, \textit{Differential geometry of lightlike submanifolds}, Frontiers in Mathematics, Birkh\"auser, Basel, 2010.
\bibitem[Furuhata2009]{Furuhata2009} H. Furuhata, Hypersurfaces in statistical manifolds, \textit{Differential Geom. Appl.} \textbf{27} (2009), no.~3, 420--429.
\bibitem[JainSinghKumar2020]{JainSinghKumar2020} V. Jain, A. P. Singh and R. Kumar, On the geometry of lightlike submanifolds of indefinite statistical manifolds, \textit{Int. J. Geom. Methods Mod. Phys.} \textbf{17} (2020), 2050099, \url{https://doi.org/10.1142/S0219887820500991}.
\bibitem[KazemiSalahvarzi2020]{KazemiSalahvarzi2020} M. B. Kazemi Balgeshir and S. Salahvarzi, Lightlike submanifolds of semi-Riemannian statistical manifolds, \textit{Balkan J. Geom. Appl.} \textbf{25} (2020), no.~2, 52--65.
\bibitem[Lauritzen1987]{Lauritzen1987} S. L. Lauritzen, Statistical manifolds, in \textit{Differential Geometry in Statistical Inference}, IMS Lecture Notes--Monograph Series, vol.~10, Institute of Mathematical Statistics, 1987, pp.~163--216.
\bibitem[RaniKaur2021]{RaniKaur2021} V. Rani and J. Kaur, Lightlike hypersurfaces of an indefinite statistical manifold, \textit{Adv. Appl. Math. Sci.} \textbf{20} (2021), no.~9, 1995--2013.
\bibitem[Vos1989]{Vos1989} P. W. Vos, Fundamental equations for statistical submanifolds with applications to the Bartlett correction, \textit{Ann. Inst. Statist. Math.} \textbf{41} (1989), no.~3, 429--450.
\end{thebibliography}
\end{document}
